\section{Redistribution: what moves between blocks}\label{sec:redistribution}

Fix an orthogonal projector $\PP$ and write $\QQ=1-\PP$. Every operator decomposes into four
blocks, and the two off-diagonal ones are the only ones that can carry anything between the
admissible subspace and its complement.

\begin{definition}[Redistribution operator]\label{def:F}
For $\Omega\in\operatorname{End}(H)$,
\begin{equation}\label{eq:F}
   \Fr(\Omega,\PP):=\PP\Omega\QQ+\QQ\Omega\PP .
\end{equation}
\end{definition}

\begin{lemma}[Block decomposition]\label{lem:blocks}
For every $\Omega$, $\ \Omega=\PP\Omega\PP+\PP\Omega\QQ+\QQ\Omega\PP+\QQ\Omega\QQ$, and the two
off-diagonal blocks of $\Fr(\Omega,\PP)$ are the off-diagonal blocks of $\Omega$:
\[
   \PP\,\Fr(\Omega,\PP)\,\QQ=\PP\Omega\QQ,\qquad
   \QQ\,\Fr(\Omega,\PP)\,\PP=\QQ\Omega\PP .
\]
\T\ Kernel-checked: \lean{block\_decomposition}, \lean{redistribution\_block12},
\lean{redistribution\_block21} (\lean{EmergentTransport/GapsAndTotals.lean}).
\end{lemma}

\begin{proof}
$\PP\Omega\PP+\PP\Omega\QQ=\PP\Omega(\PP+\QQ)=\PP\Omega$ and likewise for $\QQ$; adding gives
$\Omega$. For the blocks, $\PP\Fr\QQ=\PP(\PP\Omega\QQ+\QQ\Omega\PP)\QQ=\PP\Omega\QQ\,\QQ+\PP\QQ\,\Omega\PP\QQ=\PP\Omega\QQ$,
using $\PP\QQ=0$ and $\QQ^2=\QQ$. The other block is the same computation with $\PP$ and $\QQ$
interchanged.
\end{proof}

So $\Fr$ is exactly the off-diagonal part, and $0=\Fr$ is exactly the statement that $\Omega$ is
block-diagonal.

\begin{theorem}[Redistribution is traceless]\label{thm:traceless}
For every $\Omega$ and every idempotent $\PP$ (orthogonality is not needed here),
$\Tr\Fr(\Omega,\PP)=0$.
\T\ Kernel-checked: \lean{redistribution\_traceless}.
\end{theorem}

\begin{proof}
$\Tr(\PP\Omega\QQ)=\Tr(\Omega\QQ\PP)=0$ by cyclicity and $\QQ\PP=0$; likewise
$\Tr(\QQ\Omega\PP)=0$.
\end{proof}

\begin{theorem}[Compatibility]\label{thm:compat}
$\Fr(\Omega,\PP)=0$ if and only if $\Omega\PP=\PP\Omega$. Neither self-adjointness of $\Omega$
nor orthogonality of $\PP$ is used.
\T\ Kernel-checked: \lean{redistribution\_zero\_iff\_commute}.
\end{theorem}

\begin{proof}
($\Leftarrow$) If $\Omega\PP=\PP\Omega$ then $\PP\Omega\QQ=\PP(\Omega\QQ)=\PP(\QQ\Omega)=\PP\QQ\Omega=0$
and $\QQ\Omega\PP=(\QQ\Omega)\PP=(\Omega\QQ)\PP=\Omega(\QQ\PP)=0$.
($\Rightarrow$) If $\Fr=0$ then by Lemma~\ref{lem:blocks} both $\PP\Omega\QQ$ and $\QQ\Omega\PP$
vanish, so $\Omega=\PP\Omega\PP+\QQ\Omega\QQ$ and therefore
$\Omega\PP=\PP\Omega\PP=\PP\Omega$.
\end{proof}

The next result is the quantitative one: the \emph{size} of the movement between blocks equals
the size of the failure to commute.

\begin{theorem}[Norm identity]\label{thm:norm}
Let $\Omega$ be self-adjoint and $\PP$ an orthogonal projector. Then
\begin{equation}\label{eq:norm}
   \norm{\Fr(\Omega,\PP)}_F^2 \;=\; 2\,\norm{\PP\Omega\QQ}_F^2 \;=\; \norm{[\Omega,\PP]}_F^2 ,
   \qquad\text{and}\qquad \tfrac12\norm{[\Omega,\PP]}_F^2=\norm{\PP\Omega\QQ}_F^2 .
\end{equation}
Moreover $[\Omega,\PP]$ is skew: $[\Omega,\PP]^*=-[\Omega,\PP]$.
\T\ Kernel-checked: \lean{redistribution\_norm}, \lean{commutator\_skew},
\lean{half\_commutator\_norm}.
\end{theorem}

\begin{proof}
Put $A=\PP\Omega\QQ$ and $B=\QQ\Omega\PP$. Since $\Omega^*=\Omega$ and $\PP^*=\PP$,
$B=A^*$. Since $\QQ\PP=0$, $A^2=\PP\Omega\QQ\PP\Omega\QQ=0$, and likewise $B^2=0$.
Then $\Fr=A+B$, $\Fr^*=\Fr$, and
\[
   \norm{\Fr}_F^2=\Tr\bigl((A+A^*)^2\bigr)=2\norm{A}_F^2+2\operatorname{Re}\Tr(A^2)=2\norm{A}_F^2 .
\]
For the commutator, $[\Omega,\PP]=\Omega\PP-\PP\Omega=(\QQ\Omega\PP+\PP\Omega\PP)-(\PP\Omega\PP+\PP\Omega\QQ)=B-A=A^*-A$,
so $\norm{[\Omega,\PP]}_F^2=\Tr\bigl((A^*-A)(A-A^*)\bigr)=2\norm{A}_F^2-2\operatorname{Re}\Tr(A^2)=2\norm{A}_F^2$.
Skewness is immediate from $\Omega^*=\Omega$, $\PP^*=\PP$.
\end{proof}

\begin{remark}[A correction, recorded]\label{rem:corrected-chain}
An earlier formulation of this section chained the equalities as
$\norm{\Fr}^2=2\norm{\PP\Omega\QQ}^2=\tfrac12\norm{[\Omega,\PP]}^2$. That is false: the last
two terms differ by a factor $4$. The counterexample is $\PP=\operatorname{diag}(1,0)$,
$\Omega=\sigma_1$, for which $\norm{\Fr}_F^2=2$, $\norm{\PP\Omega\QQ}_F^2=1$ and
$\norm{[\Omega,\PP]}_F^2=2$, so $\tfrac12\norm{[\Omega,\PP]}_F^2=1\neq 2$. The correct chain is
\eqref{eq:norm}: $\Fr$ and the commutator have \emph{the same} squared norm, and the half-factor
belongs to the second identity, where it converts a squared commutator into a squared metric
block. This is the form proved in the kernel and the form used in \S\ref{sec:geometry}.
\end{remark}

\begin{example}[Non-normal $\Omega$: traceless but not balanced]\label{ex:nonnormal}
The first conclusion survives non-normality and the second does not. Take
$\Omega=\left(\begin{smallmatrix}0&1\\0&0\end{smallmatrix}\right)$, $\PP=\operatorname{diag}(1,0)$. Then
$\Fr(\Omega,\PP)=\left(\begin{smallmatrix}0&1\\0&0\end{smallmatrix}\right)$, so $\Tr\Fr=0$ and
$\norm{\Fr}_F^2=1=\norm{\PP\Omega\QQ}_F^2$; the factor $2$ of \eqref{eq:norm} is absent because
$\Omega$ is not self-adjoint ($\Omega\Omega^*\neq\Omega^*\Omega$). \V\ \lean{FalseControls/norm\_needs\_selfadjoint.lean}, which the verification requires
to \emph{fail}.
\end{example}

\begin{remark}[What the section says, in one sentence]
A transport can move structure across an interface exactly to the extent that the structural
operator fails to commute with the projector, and the amount moved is the whole of that
failure --- there is no third thing.
\end{remark}
